\part{Anforderungen, Analyse und Bewertung}
\chapterimage{Bilder_und_Co/analyse.jpg} % Chapter heading image
%----------------------------------------------------------------------------------------
% --- KAPITEL: ANFORDERUNGEN UND SZENARIEN ---
\chapter{Radarsystem Auswahl}\label{chap:szenarien}
Radarsysteme werden in der Luftfahrt an verschiedensten Orten und für verschiedenste Zwecke eingesetzt. 
Und tragen dabei essentiell zur Flugsicherheit sowohl am Boden als auch an Board bei.

\section{Ziel und Einordnung}\label{sec:szenarien-intro}
Um in dieser Arbeit etwas konkreter zu werden, wird ein Radarsystem augewählt und genauer betrachtet.
Da sonst alles nur sehr waage Umschreibungen ohne konkretes Ziel wären. 
Dieses Kapitel definiert Szenarien und Anforderungen für ein ausgewähltes Radarsystem.

\section{Verschiedene Radarsysteme}
In Flugzeugen sind meist mehrere Radare für verschieden Verwendungszwecke verbaut. Im folgenden werden 
beispielhaft verschiedene Flugzeuge betrachtet. Zu untersuchen sind die dort verwendeten Radarsysteme, 
deren Profile bzw. was man versucht zu erkennen und die entsprechenden Anforderungen. Um die Dimensionen 
zu verstehen je eine kleine Listen der Spezifikationen.

\subsection{MiG-25}
Ein Mig-25 besitzt ein Bordradar zur Luftzielsuche und Verfolgung und ein Seitensichtradar zur Aufklärung 
bzw. Bodenkartierung.
Als Bordradar wurde z.B. ein Smerch-A1 verbaut, die Spezifikationen sind: \parencite{MiGAlley_Smerch}

\paragraph{Überblick}
\begin{itemize}
  \item \textbf{Bezeichnungen:} RP-SA / RP-25 / \emph{Smerch-A} / \emph{Izdeliye 720} (MiG-25-Reihe); Ursprung aus dem verbesserten Tu-128A-Smerch.
  \item \textbf{Plattformen:} MiG-25P (später A2/A3 in Serienmaschinen); frühe Prototypen mit \emph{Smerch-A1}.
  \item \textbf{Radar-Gewicht (RP-25):} ca.\ 500~kg.
  \item \textbf{Grundwellenlänge:} \(\sim\)3\,cm; \emph{Smerch-A1} zusätzlich geheime 2\,cm-Betriebsart.
  \item \textbf{Reichweiten (A1):} Bomber-Detektion \(\approx 100\)~km; Tracking-Reichweite \(\approx 50\)~km.
  \item \textbf{ECM/Clutter:} fortlaufende Verbesserungen in Störfestigkeit und Tiefflug-Cluttertoleranz ab Serienstand.
\end{itemize} %600kw nach  wikipedia



\subsection{F-14}
Als Bordradar wurde unteranderem ein AWG-9 verwendet, die Spezifikationen sind: \parencite{Forecast1999_AWG9_APG71}
\paragraph{Abmessungen \& Masse}
\begin{itemize}
  \item \textbf{Antenne:} Slotted plane array; % Angabe übernommen wie geliefert
        \textit{14.2 cm} \,/\, \textit{36 in}\footnote{Die Antennenangabe ist in der gelieferten Quelle inkonsistent (14{,}2\,cm entsprechen nicht 36\,in). Wert wörtlich übernommen.}
  \item \textbf{Gewicht (System):} 590~kg \,/\, 1{,}300~lb
  \item \textbf{Volumen (System):} 0{,}78~m\textsuperscript{3} \,/\, 28~ft\textsuperscript{3}
\end{itemize}

\paragraph{Frequenz \& Leistung (AWG-9)}
\begin{itemize}
  \item \textbf{Frequenzband:} 8--12~GHz (X-Band)
  \item \textbf{Spitzenleistung (Peak Power):} 10~kW
  \item \textbf{Mittlere Leistung (Average Power):}
    \begin{itemize}
      \item 7~kW (Pulse-Doppler-Modus)
      \item 500~W (Pulsmodus)
    \end{itemize}
  \item \textbf{Pulsbreite:}
    \begin{itemize}
      \item Pulsmodus: 0{,}4~\(\mu\)s, 50~\(\mu\)s
      \item Pulse-Doppler: 0{,}4~\(\mu\)s, 1{,}3~\(\mu\)s, 2{,}0~\(\mu\)s, 2{,}7~\(\mu\)s
    \end{itemize}
\end{itemize}

\paragraph{Leistungsdaten}
\begin{itemize}
  \item \textbf{Max.\ Reichweite (Suche):} 213~km \,/\, 115~NM
  \item \textbf{Scanrate:} 80\(^\circ\)/s (horizontal); 2 Scans/s (vertikal)
  \item \textbf{Nachführkapazität (TWS):} bis zu 24 Ziele gleichzeitig; Raketenstarts in rascher Folge gegen bis zu 6 dieser Ziele
  \item \textbf{Zielverfolgbarer Höhenbereich:} 50~ft bis 80{,}000~ft; Zielgeschwindigkeit von niedriger Unterschall bis \(> \)~Mach~3
\end{itemize}



\subsection{A-320}
In einem A-320 sind drei verschiedene Radarsysteme verbaut: ein Wetterradar, ein Radar-Höhenmesser und 
ein Bodenradar (ATC). 
Als Wetterradar wird z.B. ein RDR-4000 verwendet, die Spezifikationen sind: \parencite{Honeywell_RDR4000_Specs}
\paragraph{Radar Processor}
\begin{itemize}
  \item \textbf{Weight:} 10.5~lbs (4.76~kg) max
  \item \textbf{Input Power:} 115~VAC (96--134~VAC), 360--800~Hz
  \item \textbf{Power Dissipated:} 150~VA nominal\\
        \hspace*{1.8em}\footnotesize(Inklusive Leistung für Transmitter/Receiver, Antennatrieb und Bedienpanel)
  \item \textbf{Environmental:} DO-160E (\(-55^\circ\)C bis \(+70^\circ\)C)
  \item \textbf{Software:} RTCA DO-178B Level~C
\end{itemize}

\paragraph{Transmitter/Receiver}
\begin{itemize}
  \item \textbf{Transmitter Method:} Pulse Compression
  \item \textbf{Weight:} 5.1~lbs (2.31~kg) max
  \item \textbf{Transmitter Frequency:} 9.375~GHz
  \item \textbf{Noise Figure:} 1.9~dB
  \item \textbf{Environmental:} DO-160E (\(-55^\circ\)C bis \(+70^\circ\)C)
  \item \textbf{Minimum Discernable Signal:} \(-124\)~dBm
\end{itemize}

\paragraph{Antenna System}
\begin{itemize}
  \item \textbf{Flat-Plate-Optionen (Gain / Beamwidth / Weight):}
\end{itemize}

\begin{center}
\begin{tabular}{lccc}
\hline
\textbf{Apertur (Durchmesser)} & \textbf{Gain (dBi, nom.)} & \textbf{3-dB-Beamwidth} & \textbf{Gewicht} \\
\hline
30\,"\  & 34.8 & 3.0\(^\circ\) & 6~lbs \\
24\,"\  & 33.0 & 4.2\(^\circ\) & 4~lbs \\
18\,"\  & 31.0 & 5.6\(^\circ\) & 3.0~lbs \\
\hline
\end{tabular}
\end{center}

\begin{itemize}
  \item \textbf{Scan Rate:} bis zu 90\(^\circ\)/s
  \item \textbf{Weight:} 16~lbs (Single), 29.5~lbs (Dual)
  \item \textbf{Environmental:} DO-160E (\(-55^\circ\)C bis \(+70^\circ\)C)
\end{itemize}

\paragraph{System Specifications}
\begin{itemize}
  \item \textbf{Max Detection Ranges:}
  \begin{itemize}
    \item 320~NM -- Weather und Ground~Map
    \item 60~NM -- Turbulence
    \item 5~NM -- Windshear
  \end{itemize}
  \item \textbf{Azimuth Coverage:}
  \begin{itemize}
    \item \(\pm 80^\circ\) -- Weather und Ground~Map
    \item \(\pm 40^\circ\) -- Windshear
  \end{itemize}
\end{itemize}

\subsection{Referenz (Automotive)}
Als vergleich werden Spezifikationen eines Bosch Frontradars angeschaut: \parencite{Bosch_FrontRadar_2025}
\begin{itemize}
  \item \textbf{Frequenzbereich:} 76 bis 77 GHz und 77 bis 81 GH
  \item \textbf{Reichweite:} bis zu 530 m
  \item \textbf{(H x B x T) Maße:} 56 mm x 76 mm x 20 mm
  \item \textbf{Verbrauch:} < 3.5 W
\end{itemize}


\section{Auswahl des Radarsystems}
Da das System beispielhaft dienen soll, sind Einsatzzweck, Ort oder Art des Systems zweitrangig. 
Die Objektklassifikation wird für diese Arbeit auf Grund dessen gewählt, dass sie verhältnismäßig verständlich 
und nachvollziehbar ist. In dieser Arbeit wird beispielhaft eine Onboard Wetterphänomänsklassifikation als 
Funktion ausgewählt. Wetterphänomäne zu identifizieren bietet hohes Potential, da auf Grund der Komplexität 
klassische Methoden schnell an ihre Grenzen geraten. Und ein Pilot kann davon profitieren, wenn er 
anstatt einer Klassifizierung zwischen Grün, Gelb und Rot zusätzlich wüsste ob es sich z.B. um 
Regen oder Hagel handelt.


\section{Szenarien beim Einsatz eines Klassifikations-Radars für Wetterphänomene}


% --- Risikobadges ---
\newcommand{\riskbadge}[2]{\colorbox{#1!15}{\textcolor{#1!70!black}{\scriptsize\textbf{#2}}}}
\newcommand{\risklow}{\riskbadge{green}{NIEDRIG}}
\newcommand{\riskmed}{\riskbadge{orange}{MITTEL}}
\newcommand{\riskhigh}{\riskbadge{red}{HOCH}}
% --- Scope-Badges (zusätzlich zu den Risk-Badges) ---
\newcommand{\scopeon}{\colorbox{cyan!15}{\textcolor{cyan!60!black}{\scriptsize\textbf{ONBOARD}}}}
\newcommand{\scopeapt}{\colorbox{violet!15}{\textcolor{violet!60!black}{\scriptsize\textbf{FLUGHAFEN}}}}
\newcommand{\scopeboth}{\colorbox{gray!20}{\textcolor{black!70}{\scriptsize\textbf{BEIDES}}}}

\subsection*{Legende}
\begin{itemize}[leftmargin=2.2em,itemsep=0.2em]
  \item \risklow\ \,geringes Fehlklassifikations-/Folgerisiko
  \item \riskmed\ \,moderates Risiko
  \item \riskhigh\ \,hohes Risiko (operativ kritisch)
  \item \scopeon\ \,relevant primär für Onboard-Wetterradar / Cockpit
  \item \scopeapt\ \,relevant primär für Flughafen-/Bodenradar
  \item \scopeboth\ \,für beide Betriebsarten relevant
\end{itemize}

\subsection*{Meteorologische Szenarien}
\begin{itemize}[leftmargin=2.2em,itemsep=0.35em]
  \item \riskhigh\ \scopeboth\ Konvektive Gewitterzellen / TS-Kerne (Thunderstorm cores) (Starkregen, Graupel, Hagel; starken Auf-/Abwinde).
  \item \riskhigh\ \scopeboth\ Eingelagerte Konvektion in Stratiformbereichen (kleinere Geweitterzellen versteckt in großflächigen harmlosen Wolken).
  \item \riskhigh\ \scopeboth\ Vereisungsrelevante Phasen: Schnee/Schneeregen (Vereisungsgefahr).
  \item \riskmed\ \scopeboth\ Graupelschauer / Wintergewitter (gemischte Wasser- und Eisteilchen, schnelle Klassendynamik).
  \item \riskmed\ \scopeboth\ Anvil-Blowoff, Virga, (Eispartikel, die von einer Gewitterwolke weggeweht wurden; Verdunstung bevor der Boden berührt wurde) (Echo ohne Bodenniederschlag; starker lokaler Abwind (Downburst) durch Verdunstungskühlung).
  \item \riskhigh\ \scopeon\ Pfadabschattung hinter Starkzellen (\emph{black hole/gap illusion}): starke Dämpfung erzeugt scheinbar „freie“ Korridore hinter einer Zelle → hohes Risiko für Fehlrouting durch vermeintliche Lücken.
\end{itemize}

\subsection*{Mess- und Auswerteartefakte}
\begin{itemize}[leftmargin=2.2em,itemsep=0.35em]
  \item \riskhigh\ \scopeboth\ Abschwächung / Dämpfung (v.\,a.\ Schatten hinter Starkzellen).
  \item \riskmed\ \scopeboth\ Strahlenaufweitung / Beam Overshooting (bodennahe Schichten/Schmelzschicht verfehlt durch Erdkrümmung; Fernbereich bzw.\ falscher Tilt).
  \item \riskmed\ \scopeboth\ Partielle Volumenfüllung (großes Strahlvolumen und kleine Zellen $\rightarrow$ wird gemittelt).
  \item \riskhigh\ \scopeapt\ Kalibration / Drift (z.\,B.\ Offset $\rightarrow$ systematische Fehler).
\end{itemize}

\subsection*{Operative Anwendungsszenarien}
\begin{itemize}[leftmargin=2.2em,itemsep=0.35em]
  \item \riskhigh\ \scopeboth\ TS-Avoidance (Klassifikation von Hagel/Starkregen; Routen-/Holding-Entscheidungen).
  \item \riskhigh\ \scopeboth\ Vereisungsmanagement (Erkennen von z.B. SLD (supercooled large droplets); kann zu schneller Leistungsdegradation führen).
  \item \riskhigh\ \scopeboth\ Mikroburst Erkennung (Kurzlebige, starke Abwinde) (start-/landekritisch).
\end{itemize}



\section{Abgeleitete Anforderungen für Onboard-/Aviation-Radar}

\subsection*{1. Gewitterkerne und Linien erkennen (\riskhigh \scopeboth)}
\begin{itemize}[leftmargin=2.2em,itemsep=0.25em]
  \item [F] Das System hebt starke Zellen und Gewitterlinien gut sichtbar hervor.
  \item [L] Anzeige aktualisiert sich schnell (ca. alle 5–10 s) und deckt mindestens 160 NM ab.
  \item [H] Klare Markierungen/Umrisse und Entfernungsangaben helfen bei Ausweichentscheidungen.
  \item [V] In Tests werden typische Gewitterszenarien zuverlässig als „zu meiden“ markiert.
\end{itemize}

\subsection*{2. Eingelagerte Zellen in großflächigem Regen (\riskhigh \scopeboth)}
\begin{itemize}[leftmargin=2.2em,itemsep=0.25em]
  \item [F] Kleine, „versteckte“ Kerne im flächigen Regen werden als eigene Gefahrbereiche angezeigt.
  \item [L] Auch bei nur wenig Helligkeitsunterschied zum Umfeld sollen diese Kerne erkannt werden.
  \item [H] Ein einfacher Hinweis („Eingelagert“) kann zugeschaltet werden.
  \item [V] Nachweis mit aufgezeichneten Fällen und Vergleichsdaten (z.\,B. Blitzmeldungen).
\end{itemize}

\subsection*{3. Vereisung erkennen und managen (\riskhigh \scopeboth)}
\begin{itemize}[leftmargin=2.2em,itemsep=0.25em]
  \item [F] Das System gibt Hinweise auf Vereisungsbereiche (inkl. feiner Niesel unter 0~\textdegree C).
  \item [L] Die ungefähre Höhe der „Null-Grad-Schicht“ wird angezeigt mit brauchbarer Genauigkeit.
  \item [H] Ein einfaches Profil entlang der Flugstrecke zeigt, wo Vereisung wahrscheinlich ist.
  \item [V] Abgleich mit Pilotenberichten und Vorhersagen zeigt eine gute Übereinstimmung.
\end{itemize}

\subsection*{4. Schauer/Graupel und schnelle Änderungen (\riskmed \scopeboth)}
\begin{itemize}[leftmargin=2.2em,itemsep=0.25em]
  \item [F] Kurze, kräftige Schauer werden erkennbar klassifiziert.
  \item [L] Starke Änderungen sollen innerhalb weniger Sekunden sichtbar werden.
  \item [H] Hinweis „schnelle Entwicklung“ zur Aufmerksamkeitserhöhung.
\end{itemize}

\subsection*{5. Amboss-Auswehung, Fallstreifen, böiger Ausfluss (\riskmed \scopeboth)}
\begin{itemize}[leftmargin=2.2em,itemsep=0.25em]
  \item [F] Bereiche mit Niederschlag in der Höhe ohne Bodentreffer sowie Böenfronten werden gekennzeichnet.
  \item [L] Die Lage solcher Grenzen wird ausreichend genau dargestellt (für taktische Meidung).
  \item [H] Kopplung mit Windscherungs-/Böenwarnungen, wenn solche Daten vorliegen.
\end{itemize}

\subsection*{6. Dämpfung und scheinbare „Lücken“ (\riskhigh \scopeon)}
\begin{itemize}[leftmargin=2.2em,itemsep=0.25em]
  \item [F] Das System erkennt und kennzeichnet abgeschattete Bereiche hinter starken Zellen.
  \item [L] Falsch zu schwach angezeigte Echos hinter einer Zelle sollen deutlich markiert werden.
  \item [H] Deutlicher Hinweis „Sicht eingeschränkt – Messung unsicher“ und Vorschlag, seitlich auszuweichen.
  \item [V] Nachweis in Flugtests, dass irreführende „Gaps“ seltener zu Fehlentscheidungen führen.
\end{itemize}

\subsection*{7. Strahlführung und „Overshooting“ (\riskmed \scopeboth)}
\begin{itemize}[leftmargin=2.2em,itemsep=0.25em]
  \item [F] Eine Automatik hilft beim richtigen Neigungswinkel; warnt, wenn tiefe Schichten übersehen werden könnten.
  \item [L] Tiefe Schichten im Nahbereich und mittlere Höhen in der Ferne sollen abgedeckt sein.
  \item [H] Einfache Empfehlung „bester Neigungswinkel“ und ein leicht bedienbarer Automatikmodus.
\end{itemize}

\subsection*{8. Ausweichen und Routenhilfe (\riskhigh \scopeboth)}
\begin{itemize}[leftmargin=2.2em,itemsep=0.25em]
  \item [F] Der geplante Flugweg wird auf Konflikte mit Zellen geprüft (einige Minuten im Voraus).
  \item [L] Vorschläge für seitlichen Abstand werden ausreichend genau berechnet.
  \item [H] Klare, einfache Empfehlungen (z.\,B. „rechts 20 NM umfliegen“) ohne Informationsüberflutung.
  \item [V] Erprobung mit Crew-Feedback und Verfahrensvorgaben.
\end{itemize}

\subsection*{9. Vereisung im laufenden Betrieb (\riskhigh \scopeboth)}
\begin{itemize}[leftmargin=2.2em,itemsep=0.25em]
  \item [F] Automatische Hinweise: „Feine Vereisung möglich“, „Eisbereich in X Minuten“.
  \item [L] Zeit- und Entfernungsangaben sind grob, aber verlässlich genug für die Planung.
  \item [H] Anbindung an Checklisten-Hinweise im Cockpit (nur als Unterstützung).
\end{itemize}

\subsection*{10. Windscherung und Fallböen (\riskhigh \scopeboth)}
\begin{itemize}[leftmargin=2.2em,itemsep=0.25em]
  \item [F] Hinweise auf starke Böengrenzen und plötzlich wechselnde Winde werden angezeigt; externe Meldungen können eingeblendet werden.
  \item [L] Warnungen erscheinen schnell genug, um darauf reagieren zu können.
  \item [H] Einheitliche, gut verständliche Warnstufen und klare Cockpit-Hinweise.
\end{itemize}

\subsection*{11. Allgemeines}
\begin{itemize}[leftmargin=2.2em,itemsep=0.25em]
  \item [L] Verfügbarkeit und Fehlermeldungen sind klar erkennbar (z.\,B. „Dämpfungskorrektur nicht aktiv“).
  \item [F] Kurze Aufzeichnung zur Nachbetrachtung (z.\,B. letzte 30 Minuten) ist möglich.
  \item [V] Entwicklung und Tests folgen den üblichen Luftfahrt-Standards; wichtige Szenarien sind durch Tests abgedeckt.
\end{itemize}


\subsection*{Legende}
\paragraph{Anforderungstypen}
\begin{itemize}[leftmargin=2.2em,itemsep=0.2em]
  \item \textbf{F = Funktionalität} \quad Was das System tun \emph{muss}.
  \item \textbf{L = Leistung/Qualität} \quad Messbare Kennzahlen/Schwellen.
  \item \textbf{H = HMI/Bedienung} \quad Anzeige, Interaktion, Workload-Reduktion.
  \item \textbf{V = Validierung/Safety} \quad Nachweis, Tests, Konformität.
\end{itemize}




% -------------------------
% NEW CHAPTER: Ausgewählte KI-Modelle
% -------------------------
\chapter{Ausgewählte KI-Modelle}\label{chap:modelle}

\section{Ziel und Einordnung}\label{sec:modelle-intro}
Dieses Kapitel dient dazu einzelne relevante KI Modelle besser zu verstehen und eines davon 
auszuwählen, um dieses im Kontext der gewählten Objektklassifikation weiter zu untersuchen.


\section{Modell-Steckbriefe}

\subsection*{CNN (Convolutional Neural Networks)}
\begin{itemize}
  \item \textbf{Zweck:} Pixel/Region-Segmentierung \& Detektion.
  \item \textbf{Typische Variante:} YOLO
  \item \textbf{Input/Features:} 2D Radardaten (z.B. Winkel, Distanz), Radarrohdaten möglich.
  \item \textbf{Stärken:} Hohe Genauigkeit; lernt komplexe Muster; edge-tauglich mit  Quantisierung.
  \item \textbf{Schwächen:} Daten-/Label-intensiv; geringere Erklärbarkeit.
  \item \textbf{Latenz/Compute:} 10--100\,ms
  \item \textbf{Modellgröße:} ca. 2.000.000 Parameter.
\end{itemize}

\subsection*{Baumbasierte Ensembles (GBDT)}
\begin{itemize}
  \item \textbf{Zweck:} Tabellarische Klassifikation auf Feature-Vektoren.
  \item \textbf{Typische Variante:} Random Forest
  \item \textbf{Input/Features:} Texturen, Spektralbreite, SNR/Attenuation/Geometrie, Kontext (Temperatur/Höhe).
  \item \textbf{Stärken:} Schnell, robust bei kleinen/mittelgroßen Datensätzen; gut erklärbar; sehr geringe Latenz.
  \item \textbf{Schwächen:} Räumliche Nachbarschaft nur indirekt; ungeeignet für rohe Bilder/Sequenzen.
  \item \textbf{Latenz/Compute:} $< 5$\,ms/Objekt; Echtzeit fähig
  \item \textbf{Modellgröße:} ca. 13.000 Knoten/Blätter.
\end{itemize}

\subsection*{Vision Transformer }

\begin{itemize}
  \item \textbf{Zweck:} Patch-/Bildbasierte Klassifikation, Segmentierung und Detektion auf Radar-„Bildern“;.
  \item \textbf{Typische Variante:} \emph{DeiT-Small} bei mehr Daten/Compute.
  \item \textbf{Input/Feature:} Mehrkanal-Radar-Slices (z.\,B.\ Reflektivität, Spektralbreite, \text{SNR}).
  \item \textbf{Stärken:} Globaler Kontext durch Self-Attention; flexible Tokenisierung (auflösungsrobust).
  \item \textbf{Schwächen:} Hoher Daten- und Compute-Bedarf; Sensitivität gegenüber Patch-/Token-Design.
  \item \textbf{Latenz/Compute:}  $\sim$10--40\,ms pro 512$\times$512 (auf einer GPU).
  \item \textbf{Modellgröße:} ca. 22.000.000 Parameter.
\end{itemize}


\subsection*{Clustering}
\begin{itemize}
  \item \textbf{Zweck:} Objektbildung, Vor-/Nachverarbeitung, Tracking.
  \item \textbf{Typische Variante:} DBSCAN
  \item \textbf{Input/Features:} (Reflektivität, Spektralbreite, Textur) + Koordinaten; für Zeitindex/Track-Kontext.
  \item \textbf{Stärken:} Formflexibel, label-arm; stabilisiert Objekte/Tracks und verbessert nachgelagerte Klassifikationen.
  \item \textbf{Schwächen:} Parameter-sensitiv ($\varepsilon$/minPts); keine semantischen Labels.
  \item \textbf{Latenz:} optional Echtzeit fähig
  \item \textbf{Modellgröße:} Feature anhängig.
\end{itemize}

\section{Modell Vergleich und Auswahl}

\small
\begin{longtable}{p{0.16\linewidth}p{0.21\linewidth}p{0.2\linewidth}p{0.21\linewidth}p{0.2\linewidth}}
\caption{Vergleich der Modelle mit Anforderungen}\\
\toprule
\textbf{Kriterium} &
\textbf{CNN} &
\textbf{Baumbasiert} &
\textbf{Probabilistisch/OOD} &
\textbf{Clustering} \\
\midrule
\endfirsthead

\toprule
\textbf{Kriterium} &
\textbf{CNN} &
\textbf{Baumbasiert} &
\textbf{Probabilistisch/OOD} &
\textbf{Clustering} \\
\midrule
\endhead

\midrule
\multicolumn{5}{r}{\emph{Fortsetzung auf der nächsten Seite}}\\
\midrule
\endfoot

\bottomrule
\endlastfoot

\textbf{Passung zu Szenarien} &
\textbf{Hoch} lernt komplexe räumliche Muster \& Kerne &
\textbf{Hoch} (mit guten Features) keine Rohdaten &
\textbf{Hoch} auf bei großräumigen Mustern &
\textbf{Niedrig als Primär} formt Objekte, aber keine Klassen \\

\textbf{Leistung (F1, $P_d$/FAR)} &
\textbf{Hoch} bei genügend Daten &
\textbf{Hoch} (mit guten Features) niedrige FAR &
\textbf{Hoch} bei genügend Daten &
\textbf{Parameterabhängig} keine Klassenmetriken \\

\textbf{Latenz} &
\textbf{Erfüllbar}&
\textbf{Sehr gut} ($<\,$ms/Objekt) &
\textbf{mittel} (GPU notwendig) &
\textbf{Gut} \\

\textbf{Datenbedarf/ Labeling} &
\textbf{Hoch}  &
\textbf{Mittel} (Objekt-Feature-Labels) &
\textbf{Hoch} &
\textbf{Gering} \\

\textbf{Erklärbarkeit} &
\textbf{Mittel}  &
\textbf{Gut}  &
\textbf{Schlecht}  &
\textbf{Niedrig} (keine Labels) \\

\textbf{Harware- anforderung} &
\textbf{Mittel}  &
\textbf{Niedrig} &
\textbf{Hoch}  &
\textbf{Niedrig} \\

\textbf{Alleinstellungs-Tauglichkeit} &
\textbf{Ja}  &
\textbf{Ja} bei guten Features &
\textbf{Ja}  &
\textbf{Nein} Vor-/Nachverarbeitung \\
\end{longtable}

Als alleinstehendes Modell empfehlen sich Baummodelle (so fern man gute Features hat), als typisches Modell wird Random Forrest verwendet. 
Diese sind alleinstellungsfähig, was für diese Arbeit wichtig ist, um zusätzliche Komplexität zu vermeiden. 
Zudem sind sie Echtzeitfähig, gut erklärbar und haben eine relativ geringe Anforderung an die Leistung der Hardware. 
Das hat den Vorteil, dass es nicht bereits an der Hardwarezulassung scheitert.

\section{Erklärung eines Random Forrest}
Ein Random Forest ist ein Ensemble vieler Entscheidungsbäume. Jeder Baum wird aus einer zufälligen Stichprobe 
der Trainingsdaten aufgebaut und betrachtet an jedem Knoten nur eine zufällige Auswahl der Merkmale. Der Baum 
stellt einfache Schwellenfragen, teilt die Daten Schritt für Schritt und gibt am Blatt eine Klassenentscheidung 
(oder einen Zahlenwert) aus. Für neue Daten werden die Ergebnisse aller Bäume zusammengeführt: Bei Klassifikation 
zählt die Mehrheit, bei Regression der Mittelwert. Durch die zufällige Daten- und Merkmalswahl werden die Bäume 
voneinander unabhängiger, die Streuung sinkt, und die Gesamtvorhersage wird robuster.

\begin{figure}[h!]
\centering
\begin{tikzpicture}[
  level 1/.style={sibling distance=46mm, level distance=14mm},
  level 2/.style={sibling distance=28mm, level distance=14mm},
  edge from parent/.style={draw,-Stealth},
  every node/.style={font=\small},
  decision/.style={rectangle,draw,rounded corners,align=center,inner sep=2pt,minimum width=22mm},
  leaf/.style={ellipse,draw,fill=gray!10,align=center,inner sep=2pt,minimum width=16mm}
]
\node[decision]{\textbf{Frage 1}\\$f_1 \le \theta_1$?}
  child { node[decision]{\textbf{Frage 2}\\$f_2 \le \theta_2$?}
    child { node[leaf]{\textbf{Blatt}\\Klasse A} edge from parent node[left]{Ja} }
    child { node[leaf]{\textbf{Blatt}\\Klasse B} edge from parent node[right]{Nein} }
    edge from parent node[left]{Ja}
  }
  child { node[decision]{\textbf{Frage 3}\\$f_3 \le \theta_3$?}
    child { node[leaf]{\textbf{Blatt}\\Klasse B} edge from parent node[left]{Ja} }
    child { node[leaf]{\textbf{Blatt}\\Klasse C} edge from parent node[right]{Nein} }
    edge from parent node[right]{Nein}
  };
\end{tikzpicture}
\caption{Einfache Skizze eines Entscheidungsbaums: An jedem Knoten wird eine Schwellenfrage gestellt; am Blatt steht die Klassenentscheidung.}
\end{figure}

\subsection{Sicherheitseinordnung}
Erfahrungsgemäß kann man bei einer eindeutigen Klassifikation ``out of the bag'' sprich ohne Aufwand 
von einer bis zu 80 \% Genauigkeit ausgehen. Mit geringen Feature Egineering bereits 90 bis 97 \%. 
Ab hier beginnt die eigentliche Arbeit. Auf Datensätzen von wenigen hundert bis mehreren Tausend Daten 
sind Genauigkeiten bis zu 99,9 \% möglich. Mit größeren Datenmengen, guten Features und eng gefasster 
ODD wäre somit eine Genauigkeit im Bereich $\le 1\times10^{-4}$ bzw. >99,99 \% auf homogenen Testdaten realistisch. 
Da Daten in der Realität nicht homogen sondern heterogen verteilt sind, wäre eine Fehlerquote von $\ 1\times10^{-6}$ 
pro Mission theoretisch möglich. Dadurch könnten unter idealen Vorraussetzungen und besonderer Begründung 
Bereiche wie DAL B oder ASIL C möglich sein. Wenn man keine eindeutige Klassifikation sondern 
prozentuale Ausgaben berücksichtigt kann man zudem unterstützende Einsatzmöglichkeiten erörtern.

\subsection{Wichtige Punkte}

\paragraph{Eigenschaften}
\begin{itemize}
  \item Ensemble aus vielen Entscheidungsbäumen mit Zufall in Daten- und Merkmalsauswahl (\emph{robust, stabil}).
  \item Trifft Entscheidungen über einfache Schwellenfragen; am Blatt Mehrheitsklasse bzw.\ Mittelwert.
  \item Funktioniert ohne starke Annahmen über Datenverteilungen; unempfindlich gegenüber Ausreißern und Skalierung.
  \item Liefert Klassenwahrscheinlichkeiten (als Häufigkeiten im Blatt) und Feature-Wichtigkeiten.
\end{itemize}

\paragraph{Zentrale Hyperparameter (Klassifikation)}
\begin{itemize}
  \item \texttt{n\_estimators}: Anzahl Bäume (typ.\ 50--300).
  \item \texttt{max\_depth}, \texttt{min\_samples\_leaf}: Baumkomplexität/Glättung (z.\,B.\ Tiefe 6--12; Leaf $\geq$ 5--50).
  \item \texttt{max\_features}: Anzahl Merkmale pro Split (z.\,B.\ $\sqrt{d}$).
  \item \texttt{class\_weight}: Ausgleich unausgewogener Klassen (\texttt{balanced}).
  \item \texttt{criterion}: Gütemaß für Splits (Gini/Entropie).
\end{itemize}

\paragraph{Vorteile}
\begin{itemize}
  \item Geringes Overfitting-Risiko im Vergleich zu einzelnen tiefen Bäumen; solide Basisleistung mit wenig Tuning.
  \item Sehr schnelle Inferenz auf Objekt-/Patch-Features (ms-Bereich); gut für Edge/Onboard.
  \item Erklärbar über Feature-Wichtigkeiten.
\end{itemize}

\paragraph{Nachteile}
\begin{itemize}
  \item Nutzt räumliche Nachbarschaft nur indirekt; für reine Pixel-Segmentierung weniger geeignet.
  \item Empfindlich bei vielen irrelevanten Features: Zu tiefe Bäume können Rauschen „aufsammeln“.
  \item Klassengewichte/Imbalance: Ohne passende class-weight/Sampling-Strategien werden Minderheitsklassen leicht benachteiligt.
\end{itemize}


\section{Einbindungsmöglichkeiten}
Ein Random Forrest kann an verschiedenen Stellen eingesetzt werden, mögliche Einsatzfelder sind:

\begin{itemize}
  \item \textbf{Musteridentifikation \& Datendarstellung in der Entwicklung}\\
  Wiederkehrende Strukturen in Radardaten können mit einem Random Forrest systematisch erkannt und als 
  verständliche Merkmalsbilder aufbereitet werden. Hierzu werden geeignete Features (z.\,B.\ Intensitäts-, Textur- und Kontextgrößen) definiert, 
  Daten explorativ visualisiert und ein konsistenter „Musterkatalog“ erstellt, der später als Grundlage für die Klassifikation dient.

  \item \textbf{Assistierendes System neben dem klassischen Radar}\\
  Das Modell läuft parallel und liefert zusätzliche, verdichtete Informationen zur Wetterlage 
  (z.\,B.\ Klassenhinweise, Trendindizien, Qualitätsmerkmale). Die Ausgabe ergänzt die bestehende 
  Anzeige, ohne deren Logik zu ersetzen, und unterstützt die Crew mit klaren, kompakten Hinweisen 
  für die taktische Entscheidungsfindung.

  \item \textbf{Ersatz/Ergänzung klassischer Schwellen-Entscheidungsbäume}\\
  Anstelle fest verdrahteter Schwellwerte werden datengetriebene Entscheidungsgrenzen verwendet, 
  die sich aus gelabelten Beispielen ableiten lassen. Das reduziert Hand-Tuning, erhöht die Konsistenz 
  über verschiedene Wetterregime und ermöglicht feinere Trennungen zwischen ähnlichen Phänomenen.

  \item \textbf{Zusammen mit Kalman-Filter}\\
  Das Modell ersetzt die Zustands­schätzung klassischer Filter nicht, sondern \emph{füttert} sie: 
  Es liefert stabilisierte Beobachtungen und Klassifikationsergebnisse (inkl.\ Qualitätsmerkmalen), 
  die anschließend von einem Schätzer weiterverarbeitet werden. So bleibt die Bahn-/Zellverfolgung im 
  Filter, während die inhaltliche Deutung der Echos und ihre zeitliche Glättung modellbasiert verbessert werden.
  Der Kalmanfilter (ähnlich eines hidden Markow Modells) wurde als Ansatz ausgesucht, da dessen Output ebenfalls 
  probabilistisch ist und daher durch die Ergänzung eines ebenfalls probabilistischen Algorithmus nicht 
  zwangsläufig schwerwiegend beeinflusst wird.
\end{itemize}

Die Musteridentifikation und Datendarstellung in der Entwicklung bietet das größte Potenzial, weil sie der 
EASA-Level-1-Kategorie zugeordnet und damit zeitnah eingesetzt werden kann. Zudem ist durch die Arbeit mit 
großen Datensätzen zu erkennen, wie wertvoll ML-Verfahren sind. Das zeigt, dass die Entwicklung von der Disziplin Data Science 
sehr profitieren kann.

\section{Mögliche Szenarien beim Einsatz}

\begin{table}[h!]
\centering
\small
\begin{tabular}{p{0.21\linewidth}p{0.37\linewidth}p{0.38\linewidth}}
\toprule
\textbf{Szenario} & \textbf{Beschreibung} & \textbf{Erklärung (Bezug zu klassischen Methoden / bekanntem Wissen)} \\
\midrule
Übereinstimmung mit klassischer Methode & Random Forest und klassische Auswertung kommen zur gleichen Aussage. & Es werden etablierte Verfahren bzw.\ Vorwissen bestätigt; kein besonderer Erkenntnisgewinn. \\
\addlinespace
Nahe beieinander / knapp verschieden & Random Forest und klassische Methode liefern ähnliche, aber nicht identische Ergebnisse. & Grenzbereich: Daten/Features geben keine klare Trennung her oder Kontexteffekte fehlen; Entscheidungen können schwanken und sollten markiert bzw.\ weiter untersucht werden. Hinweis, dass Daten nicht vollständig oder repräsentativ waren, dass das Modell angepasst werden muss oder vielleicht bekanntes Wissen detaillierter werden kann.\\
\addlinespace
Einseitige Lücke & Entweder die klassische Methode deckt einen Fall nicht ab, den der Random Forest sinnvoll bewertet – \emph{oder} der Random Forest versagt, wo die klassische Methode stabil ist. & Hinweis auf Abdeckungsunterschiede: z.\,B.\ klassische Schwellen nicht definiert für neuen Kontext, oder das Modell außerhalb seiner Trainingsdomäne; klare Dokumentation, wann welches Verfahren gilt. \\
\addlinespace
Gegenteilige Aussagen & Random Forest und klassische Methode widersprechen sich deutlich. & Konfliktfall: mögliche Label-/Feature-Probleme, Hyperparameter Wahl, etc.; erfordert gezielte Fallanalyse. \\
\addlinespace
Neues Muster entdeckt & Random Forest macht auf wiederkehrende Strukturen aufmerksam, die im bisherigen Wissen/Katalog nicht vorkommen. & Kandidat zur Erweiterung des Musterkatalogs: Fälle sammeln, prüfen, definieren und (falls valide) als neue Klasse/Regel aufnehmen. Hier gilt ``kein Assurance Rabatt'', es muss mit etablierten Methoden nachgewießen werden. \\
\bottomrule
\end{tabular}
\caption{Szenarien bei der Musteridentifikation \& Datendarstellung mit Random Forest im Vergleich zu klassischen Methoden bzw.\ bekanntem Wissen.}
\end{table}


% -------------------------
% Anforderungsfelder & Referenzrahmen
% -------------------------
\chapter{Anforderungsfelder \& Referenzrahmen}\label{chap:anforderungen}

\section{Ziel des Kapitels}\label{sec:anforderungen-intro}
Dieses Kapitel fasst wiederkehrende Sicherheitsbedarfe als \emph{Anforderungsfelder} zusammen und ordnet ihnen
\emph{typische Erwartungen} aus etablierten Leitlinien und Standards zu.
%die Felder bilden die Zeilen der Coverage-Analyse in Kapitel~\ref{chap:coverage}. 

\section{Anforderungsfelder}\label{sec:anforderungsfelder}

\subsection*{1. Klassentaxonomie \& Entscheidungslogik}
\begin{itemize}[leftmargin=2.2em,itemsep=0.2em]
  \item Eindeutige, domänenspezifische Klassen mit Kurzdefinition und Abgrenzung (z.\,B. TS-Kern vs. starker Regen; Graupel vs. Hagel).
  \item Qualitative Entscheidungsmerkmale je Klasse (z.\,B. Intensität, Gradienten, Persistenz) und typische Verwechslungen.
  \item \emph{Granularität/Hierarchie:} Aggregatklassen für grobe Analysen vs. Feinklassen für Detailstudien.
\end{itemize}

\subsection*{2. Semantik, Kontext \& Label-Policy}
\begin{itemize}
  \item Raum-/Zeit-Kontext: Gültigkeitsbereich der Klassendefinition (Größe, Dauer, Umgebung).
  \item Quellenbezug: Onboard-Only vs. Fusion (Bodenradar/Blitz/NWP) — was darf zur Entscheidung beitragen?
  \item Regeln für Mischpixel/Mischzellen (Primärklasse, Mehrfachlabel, „Gemischt“) und Konfliktlösung (Tie-Break, „unentschieden“).
  \item Dokumentierte Edge-Cases (z.\,B. Virga über trockener Schicht, Attenuationslücke).
\end{itemize}

\subsection*{3. Unsicherheit, OOD}
\begin{itemize}
  \item Gültige Ausgänge: \emph{Klasse A}, \emph{Klasse B}, \emph{Unbekannt/prüfen}; kein Zwang zur Entscheidung.
  \item Definition von OOD (außerhalb Trainingsspektrum; z.\,B. Chaff, seltene Lagen) und Umgang: \emph{Unbekannt} + Hinweis.
  \item Pflege eines Katalogs seltener Phänomene (Wissensbasis) zur späteren ODD-Erweiterung.
\end{itemize}

\subsection*{4. Reproduzierbarkeit}
\begin{itemize}
  \item Fixierte Datenstände, Seeds, Preprocessing-Versionen (gleicher Input \textrightarrow{} gleicher Output).
  \item Wiederholbare Experimente (Skripte/Notebooks, protokollierte Konfigurationen).
  \item Versionierung von Daten, Features, Modellen und Label-Guidelines; optional Re-Training/Updates dokumentieren.
\end{itemize}

\subsection*{5. Validierungslogik}
\begin{itemize}
  \item Strikte Trennung von Trainings- und Prüfdatensätzen (nach Ereignis/Zeitraum/Region).
  \item Klassenbezogene, qualitative Erfolgskriterien (z.\,B. „Kernbereiche werden klar isoliert“).
  \item Bedeutungsstabile Ergebnisse über Distanzen/Tilts/Quellen (sofern genutzt).
\end{itemize}

\subsection*{6. ODD, Datenraum \& Szenarioabdeckung}
\begin{itemize}
  \item ODD klar beschreiben (Regionen, Jahreszeiten, Entfernungen/Tilts, verfügbare Kanäle, Flugphasen).
  \item Szenariokatalog führen (typisch \& Edge Cases), Coverage dokumentieren, Lücken benennen.
  \item Grenzen \& Annahmen explizit (z.\,B. „keine Polarimetrie“, „Küstenclutter“); pro ODD-Segment Leistung \emph{und} Unsicherheit berichten.
  \item Änderungen am ODD/Daten/Pipeline \textrightarrow{} erneute Bewertung inkl. Regressionstestplan; kurze Kriterien, wann ein neues ODD-Segment als „abgedeckt“ gilt.
\end{itemize}

\subsection*{7. Erklärbarkeit (Explainability)}
\begin{itemize}
  \item Global: Welche Merkmale sind \emph{typisch} entscheidend je Klasse (Feature-Rangfolgen, Regelbeispiele)?
  \item Lokal: Kurze Fall-Erklärung „Warum wurde hier Klasse X gewählt?“ (Top-Merkmale, Nachbarschaft).
  \item Konsistenz: Erklärungen müssen zur Klassen-Definition passen (kein „Widerspruch zur Semantik“).
\end{itemize}



\section{Referenzrahmen: Wie die Quellen zusammenspielen}\label{sec:ref-standards}

% ===== Literatur-/Standards-Factsheets =====
\subsection{Factsheets: Leitlinien, Standards und ihre Rolle}
Die folgende Übersicht konsolidiert zentrale Quellen (Rolle im Rahmen, Stärken/Grenzen, Nutzung in dieser Arbeit) in
komprimierter Form. Sie dient als Nachschlage- und Verweisgrundlage, um Wiederholungen im Text zu vermeiden.
\emph{Wichtig:} Die Tabelle stellt \textbf{keine Bewertung oder Rangfolge} dar, sondern eine neutrale Einordnung mit
klarer Zweckbindung für diese Arbeit.
% --- optionales Feintuning (kann in die Präambel) ---
%\usepackage{longtable,booktabs,array,ragged2e,enumitem}
\newcolumntype{P}[1]{>{\RaggedRight\arraybackslash}p{#1}}
\renewcommand{\arraystretch}{1.08}
\setlength{\tabcolsep}{4.5pt}

% =======================
% Factsheets als LONGTABLE
% =======================
\begin{longtable}{P{0.17\textwidth} P{0.25\textwidth} P{0.38\textwidth} P{0.16\textwidth}}
\caption{Factsheets: Leitlinien/Standards – Rolle, Stärken/Grenzen, Nutzung}\label{tab:factsheets}\\
\toprule
\textbf{Quelle} & \textbf{Rolle im Rahmen} & \textbf{Stärken / Grenzen} & \textbf{Bemerkung} \\
\midrule
\endfirsthead

\toprule
\textbf{Quelle} & \textbf{Rolle im Rahmen} & \textbf{Stärken / Grenzen} & \textbf{Bemerkung} \\
\midrule
\endhead

\addlinespace[4pt]
\midrule
\multicolumn{4}{r}{\emph{Fortsetzung auf der nächsten Seite}}\\
\endfoot

\addlinespace[4pt]
\bottomrule
\endlastfoot

\textbf{ISO/IEC TR 5469\parencite{iso-iec-tr-5469-2024}} &
Domänenneutrale Brücke für \emph{Assurance} von KI; Prozess-Rahmen neben DO-178C. &
\begin{minipage}[t]{\linewidth}\vspace{2pt}
\begin{itemize}[leftmargin=*,itemsep=1pt,topsep=0pt]
  \item \textbf{+} Anschlussfähige Argumentationsstruktur, domain-agnostisch
  \item \textbf{–} Sehr high-level; ohne Domänenkonkretisierung nicht direkt anwendbar
\end{itemize}
\vspace{-4pt}\end{minipage}
&
Vokabular \& Struktur für die Coverage-Argumentation. \\

\textbf{ISO/PAS 8800\parencite{iso-pas-8800-2024}} &
KI-bezogene Safety-Erwartungen (Automotive) als Orientierungsrahmen. &
\begin{minipage}[t]{\linewidth}\vspace{2pt}
\begin{itemize}[leftmargin=*,itemsep=1pt,topsep=0pt]
  \item \textbf{+} Konkrete Felder (ODD, Daten, Robustheit, Unsicherheit, Betrieb)
  \item \textbf{–} Automotive-Kontext; Zulassungswege abweichend von Luftfahrt
\end{itemize}
\vspace{-4pt}\end{minipage}
&
Mapping auf Anforderungsfelder \& Coverage-Matrix. \\

\textbf{EASA AI Concept Paper (Issue 2, 2024)\parencite{easa-ai-concept-issue2-2024}} &
Behördliche Erwartungshaltung für L1/L2; Ergänzung zum bestehenden Safety-Regelwerk. &
\begin{minipage}[t]{\linewidth}\vspace{2pt}
\begin{itemize}[leftmargin=*,itemsep=1pt,topsep=0pt]
  \item \textbf{+} Klare Prinzipien (ODD, Learning Assurance, Monitoring, XAI)
  \item \textbf{–} Keine normativen Detailanforderungen; Auslegungsspielräume
\end{itemize}
\vspace{-4pt}\end{minipage}
&
Regulatorischer Kontext \& Leitplanken für Evidenzformen. \\

\textbf{DO-178C / DO-330\parencite{do-178c,do-330}} &
Software-Assurance (DO-178C, DAL-skaliert); Tool-Qualifizierung (DO-330). &
\begin{minipage}[t]{\linewidth}\vspace{2pt}
\begin{itemize}[leftmargin=*,itemsep=1pt,topsep=0pt]
  \item \textbf{+} Etablierte Nachweiswege, starke Traceability/Unabhängigkeit
  \item \textbf{–} Kein KI-Fokus (Daten/Modelle/ODD fehlen); ML-Toolchains herausfordernd
\end{itemize}
\vspace{-4pt}\end{minipage}
&
Baseline für Airborne-Software; Querverweis Toolchain/Traceability. \\

\textbf{ISO 21448 (SOTIF)\parencite{iso-21448-sotif}} &
Sicherheit der beabsichtigten Funktion; Fokus auf Wahrnehmungs-/Szenario-Lücken. &
\begin{minipage}[t]{\linewidth}\vspace{2pt}
\begin{itemize}[leftmargin=*,itemsep=1pt,topsep=0pt]
  \item \textbf{+} Szenario-/Abdeckungslogik sehr anschlussfähig
  \item \textbf{–} Automotive; keine luftfahrtspezifischen DAL/Prozesse
\end{itemize}
\vspace{-4pt}\end{minipage}
&
ODD-/Szenariologie als Referenz für Test \& Reporting. \\

\textbf{ISO 26262 (Functional Safety)\parencite{iso-26262}} &
Road-vehicles-Standard für funktionale Sicherheit; Prozessdisziplin (V-Modell), Unabhängigkeit, ASIL-Skalierung. &
\begin{minipage}[t]{\linewidth}\vspace{2pt}
\begin{itemize}[leftmargin=*,itemsep=1pt,topsep=0pt]
  \item \textbf{+} Klare Rollen/Unabhängigkeit; hazard-basierte Skalierung (ASIL)
  \item \textbf{+} Durchgängige Traceability \& Lebenszyklusdisziplin
  \item \textbf{–} Automotive-Kontext; keine DAL/DO-Bezüge, kein KI-spezifischer Fokus
\end{itemize}
\vspace{-4pt}\end{minipage}
&
Vergleichsfolie für Traceability, Rollen \& Lifecycle; Einordnung Toolchain/Änderungen. \\

\textbf{ISO/IEC 22989 (AI Begriffe)\parencite{iso-iec-22989}} &
Domänenneutrale KI-Terminologie (Model, Dataset, Bias, Uncertainty etc.). &
\begin{minipage}[t]{\linewidth}\vspace{2pt}
\begin{itemize}[leftmargin=*,itemsep=1pt,topsep=0pt]
  \item \textbf{+} Einheitliche Begriffsverwendung; dient als Nachschlagewerk
  \item \textbf{–} Definitionsstandard; keine Anforderungen/Nachweiswege
\end{itemize}
\vspace{-4pt}\end{minipage}
&
Begriffsreferenz (z.\,B. ODD, Unsicherheit, Kalibrierung). \\

\textbf{EASA AI Roadmap 2.0 (2023)\parencite{easa-ai-roadmap-2-2023}} &
Strategischer Rahmen (human-centric KI; Level-Pfad; Lebenszyklus-Blick). &
\begin{minipage}[t]{\linewidth}\vspace{2pt}
\begin{itemize}[leftmargin=*,itemsep=1pt,topsep=0pt]
  \item \textbf{+} Stufenpfad \& Vertrauensbausteine (Learning Assurance, Oversight)
  \item \textbf{–} High-level; keine normativen Detailanforderungen/DAL-Ableitung
\end{itemize}
\vspace{-4pt}\end{minipage}
&
Verortung der Assurance-Bausteine (ODD, Monitoring, Learning Assurance). \\

\textbf{EASA MLEAP Final Report (2024)\parencite{easa-mleap-final-2024}} &
Behördennahes Forschungsprojekt zu \emph{Learning Assurance} (W-Shape, Use-Cases). &
\begin{minipage}[t]{\linewidth}\vspace{2pt}
\begin{itemize}[leftmargin=*,itemsep=1pt,topsep=0pt]
  \item \textbf{+} Konkrete Methoden \& Pipeline; Evidenzen aus Luftfahrt-Use-Cases
  \item \textbf{–} Guidance/Empfehlung (nicht bindend); Fokus auf Level-1/2
\end{itemize}
\vspace{-4pt}\end{minipage}
&
Referenz für Evidenzformen (segmentiertes Reporting, Kalibrierung, ODD, Degradation). \\

\end{longtable}

\noindent\textit{Hinweis.}
ISO/PAS~8800 adressiert weitgehend dieselben Assurance-Bausteine (ODD/Annahmen, Daten-/Label-Governance,
Robustheit/OOD, Kalibrierung, Monitoring/Degradation) wie das \emph{EASA AI Concept Paper, Issue~2};
ISO/IEC~TR~5469 dient dabei als domänenneutrale Brücke für Terminologie und Argumentationsstruktur.
\parencite{iso-pas-8800-2024,easa-ai-concept-issue2-2024,iso-iec-tr-5469-2024}



\section{Luftfahrt-Zertifizierung und Forschung — Ergänzungen}

\subsection*{\textit{ML meets aerospace: challenges of certifying airborne AI} (Frontiers, 2024)\parencite{ml-meets-aerospace-2024}}
  \textbf{Fokus:} Hürden bei der Zertifizierung luftfahrtnaher KI (ODD-Disziplin, OOD/Abstain, V\&V-Lücken).\\
  \textbf{Beitrag:} Untermauert die Verschiebung von „Accuracy“ zu \emph{Assurance}; stützt Felder wie ODD, Robustheit, Kalibrierung, Toolchain.\\
  \textbf{Grenzen:} Überblicksarbeit; keine radar-spezifischen KPIs/Zulassungsfälle.

\subsection*{\textit{Towards certifiable AI in aviation: landscape, challenges, and opportunities} (2024)\parencite{towards-certifiable-ai-aviation-2024}}
  \textbf{Fokus:} Landscape/State of the Art zu Zertifizierungslandschaft, Evidenzbausteinen und offenen Punkten.\\
  \textbf{Beitrag:} Praktische Strukturierungshilfe/Checkliste für Evidenzpakete und die Coverage-Matrix in dieser Arbeit.\\
  \textbf{Grenzen:} Überwiegend konzeptionell; empirische Nachweise nur punktuell.







%---------------------------------------------
%noch komplett unsauber
%---------------------------------------------

\chapter{Luftfahrtrahmenbedingungen}

% ==== Luftfahrt ====
\section{Ziel, Motivation und Scope (Luftfahrt)}
\paragraph{Ziel.}
Einordnung, wie „sicher genug“ für KI-Funktionen in der Luftfahrt regulatorisch begründet und nachgewiesen wird: Rolle von EASA (AI Roadmap, Concept Paper Issue~2), MLEAP-Ergebnissen und etablierten Standards (z.B. DO-178C).

\paragraph{Motivation.}
KI hält in sicherheitsrelevanten Avionikfunktionen Einzug (z.\,B.\ Radar-Objektklassifizierung). Für den Einsatz im Flugbetrieb braucht es eine belastbare Begründung „sicher genug“ – auf Basis der etablierten Luftfahrt-Standards, ergänzt um KI-spezifische Evidenzen. Gleichzeitig bleiben zentrale Fragen offen (Generalisierung und Drift, Daten-/Modell-Governance, Updates im Betrieb, OOD/Abstain). Ein kurzer Seitenblick auf militärische Anwendungen ist hilfreich, weil dort Degradation, Resilienz gegen Störungen (Jamming/Täuschung) und konsequente Human-in-the-Loop-Konzepte geschärft sind; diese Aspekte liefern Anregungen, stehen hier aber nicht im Mittelpunkt der zivilen Zulassung.

\paragraph{Scope.}
Fokus auf Safety-Aspekte von KI in der Avionik, insbesondere \emph{Radar-Objektklassifizierung} als Beispielanwendung. Abdeckung der EASA-Guidance zu Level~1/2 (Assistenz, Human-AI-Teaming), \emph{offline} trainierten, zur Freigabe eingefrorenen Modellen.





\section{EASA AI-Roadmap – aktueller Stand}
Die EASA positioniert KI als „human-centric“ und staffelt Anwendungen nach Autorität: \textbf{Level 1} (Assistenz), 
\textbf{Level 2} (Human-AI-Teaming) und später \textbf{Level 3} (fortgeschrittene Automatisierung). 
Roadmap 2.0 definiert Vertrauenswürdigkeits-Bausteine (u.\,a. Learning Assurance, Erklärbarkeit, Human Oversight) 
und skizziert einen gestuften Zulassungspfad. \emph{Kernbotschaft:} KI ergänzt, aber ersetzt das bestehende Safety-Regelwerk nicht.
\paragraph{Was das praktisch bedeutet.}
\begin{itemize}
  \item \textbf{Stufenweise Zulassung:} Erst Level-1, später Level-2 – jeweils mit klar definiertem Einsatzbereich und menschlicher Aufsicht. \emph{Kein „Assurance-Rabatt“:} Die Strenge der Nachweise richtet sich weiterhin nach Hazard/DAL.
  \item \textbf{Zusatz statt Ersatz:} KI-spezifische Evidenzen (z.\,B. Daten-/Modell-Governance, Kalibrierung, OOD/Abstain, Monitoring) kommen \emph{on top} zur etablierten Software-/System-Assurance.
  \item \textbf{Lebenszyklus im Blick:} Betriebliches Monitoring, kontrollierte Updates und nachvollziehbare Änderungen sind Teil des Sicherheitsarguments („living“ Safety Case).
\end{itemize}


\section{EASA AI-Guidance 2024 (Concept Paper Issue 2) – Level 1 \& 2}
Das \emph{Concept Paper Issue 2} (03/2024) beschreibt die schrittweise Einführung von KI in sicherheitsrelevanten Anwendungen: zuerst \textbf{Level 1} (Assistenzfunktionen), anschließend ausgewählte \textbf{Level 2}-Anwendungen (\emph{Human–AI Teaming} unter klarer menschlicher Aufsicht). Die Guidance ergänzt das bestehende Safety-Regelwerk (z.\,B.\ DO-178C) um KI-spezifische Bausteine wie \emph{Learning Assurance}, Erklärbarkeit, präzise Einsatzgrenzen (ODD) und betriebliches Monitoring – ohne die klassischen Nachweise zu ersetzen.

\paragraph{Schwerpunkte.}
\begin{itemize}
  \item \textbf{Klare Rollenverteilung:} L1 = Unterstützung; L2 = Kooperation mit dem Menschen. Autorität, Eingriffsrechte und Übersteuerbarkeit müssen beschrieben und gegeben sein.
  \item \textbf{Einsatzbereich festlegen (ODD):} Scope, Use Cases, Datenbasis und Grenzen definieren – inkl.\ Regeln, wann „Abstain/Unknown“ zu ziehen ist.
  \item \textbf{Learning Assurance:} Herkunft/Qualität der Trainingsdaten, kontrollierte Modelländerungen und Freigabeprozesse transparent darlegen.
  \item \textbf{Erklärbarkeit \& HMI:} Begründbare Ausgaben, Konfidenz-/Qualitätsanzeigen und eindeutige Degradationsmodi.
  \item \textbf{Betrieb \& Monitoring:} Drift-Überwachung, Logging, Eskalationswege; \emph{„living“ Safety Case} statt Einmal-Nachweis.
\end{itemize}

\paragraph{Frühe Einsatzgrenzen.}
\begin{itemize}
  \item \textbf{Begrenzte Autorität:} L1 = Mensch entscheidet und handelt; L2 = System darf handeln, bleibt aber jederzeit übersteuerbar.
  \item \textbf{Kein Online-Lernen:} Zulässig sind \emph{offline} gelernte, zur Freigabe \emph{eingefrorene} Modelle; Online-/Continual Learning ist (noch) ausgeschlossen.
  \item \textbf{ODD \& OOD:} Der Einsatzbereich ist eng zu beschreiben; für \emph{außerhalb der ODD}/\emph{Out-of-Distribution} sind Monitoring, Warnung und sicheres Verhalten gefordert.
  \item \textbf{Nachweisstrenge bleibt:} KI-spezifische Evidenzen kommen \emph{on top}; die Software-/Hardware-Assurance (z.\,B.\ DO-178C/DO-254) bleibt vollumfänglich anwendbar – es gibt keine Erleichterung der Nachweisstrenge.
  \item \textbf{Zielniveau:} Mindestens das Sicherheitsniveau konventioneller Lösungen.
\end{itemize}




\section{MLEAP Final Report (2024) – Kernaussagen und Relevanz}
Der \emph{Final Report} des EASA-Forschungsprojekts \textbf{MLEAP} (Machine Learning Application Approval) bündelt Methoden- und Werkzeugempfehlungen zur \emph{Learning Assurance} und zeigt, wie der EASA-\emph{W-Shape} (lernprozessspezifische Ergänzung zum V-Modell) praktisch umgesetzt werden kann. Schwerpunkte sind \textbf{Datenrepräsentativität}, \textbf{Generalisierung} und \textbf{Robustheit/Stabilität} von ML-Modellen; ergänzt wird dies durch eine generische Entwicklungs-Pipeline im Sinne des W-Shape. Die Ergebnisse fließen in das \emph{EASA AI Concept Paper Issue 2} ein.

\paragraph{Kerninhalte.}
\begin{itemize}
  \item \textbf{W-Shape als Rahmen:} Abgrenzung System-Ebene vs.\ KI-Ebene; spezifische Ziele und Verifikationen für Datenmanagement, Lernprozess, Modellverifikation und Integration werden entlang des Lernzyklus verankert. 
  \item \textbf{Drei technische Achsen:} (1) \emph{Datenvollständigkeit \& Repräsentativität} inkl.\ Corner-/Edge-Cases und Auswahlraster für Datenqualifikationsmethoden; (2) \emph{Generalisierung} (Vermeidung von Unter/Überanpassung, Abschätzung/Bounds, prozessnahe Maßnahmen); (3) \emph{Robustheit \& Stabilität} gegenüber Störungen/Rauschen inkl.\ Ansätzen von empirischen bis formalen Verfahren. 
  \item \textbf{Generische Pipeline:} Ein praktikabler Ablauf zur Umsetzung des W-Shape mit klaren Prüfbausteinen (Daten-, Lern-, Performance- und Integrationsverifikationen) wird vorgeschlagen. 
  \item \textbf{Use-Cases (Evidenz):} \emph{ATC-Speech-to-Text}, \emph{Automated Visual Inspection} (In-Service-Schadenerkennung) und \emph{ACAS Xu} (UAS-Air-to-Air-Avoidance) wurden untersucht, um Methoden und Empfehlungen an realen Luftfahrtaufgaben zu validieren. 
\end{itemize}





\section{DO-178C – Rolle im KI-Kontext}
\textbf{DO-178C} bleibt die \emph{Baseline} für Software-Assurance in der Avionik: Ziele und Evidenzen skalieren mit 
\textbf{DAL A–E}. EASA verweist explizit auf die DO-178C-Daten/Liefergegenstände (inkl. Tool-Qualifikation) als Nachweisgrundlage; 
KI-spezifische Guidance kommt \emph{zusätzlich} obenauf (z.\,B. Learning Assurance), ändert aber die Software-Pflichten nicht. 
\emph{Praxis:} Daten-/Modell-Pipelines berühren Tool-Qualifikation (DO-330).

\paragraph{Bezug zu KI Anwendungen und EASA}
\begin{itemize}
  \item \textbf{Airborne Software bleibt DO-178C-pflichtig:} Inferenz-Software, Laufzeit/Framework, Treiber und Integrationscode (inkl.\ Vor-/Nachverarbeitung) werden wie gewohnt gemäß DO-178C nachgewiesen; die KI-Guidance kommt \emph{on top} (kein „DAL-Rabatt“).
  \item \textbf{Modelle sind konfigurationsgeführte Artefakte:} Gewichte/Graph/Quantisierung gelten als freigegebener Software-Stand („gefroren“ zur Zulassung) und werden wie Software-Daten behandelt; Updates sind kontrollierte Änderungen, kein Online-Lernen.
  \item \textbf{Tool-Qualifikation denken:} Trainings-/Konvertierungs-/Code-Gen-Tools können qualifikationspflichtig sein (DO-330), wenn ihre Fehler unentdeckt ins Produkt wirken könnten.
  \item \textbf{Traceability bleibt zentral:} Anforderungen \(\rightarrow\) Architektur/Code \(\rightarrow\) Tests/Analysen müssen lückenlos verknüpft sein; KI-spezifische Evidenz (z.\,B. Daten-/Modell-Governance, Kalibrierung, OOD/Abstain) wird \emph{zusätzlich} dokumentiert.
\end{itemize}
\paragraph{Wichtig.}
DO-178C und DO-330 beziehen sich nicht spezifisch auf KI, sondern allgemein auf Software und Tools. Deterministischer Code 
und MC/DC-Coverage sind stark empfohlen und gelten als Best Practice; bei KI-Komponenten sind diese Punkte teils nicht erfüllbar 
und werden in den Normen nicht ausdrücklich gefordert. Fazit: KI ist unter diesen Standards grundsätzlich zulassungsfähig, jedoch 
steigen die Anforderungen an die Nachweisführung.






%----------------------------------------------------
%----------------------------------------------------





\section{Brücke zur Coverage-Analyse}\label{sec:anforderungen-bruecke}
Die in diesem Kapitel definierten Anforderungsfelder und ihre typischen Erwartungen werden in Kapitel~\ref{chap:coverage}
zu einer \emph{Coverage-Analyse} verdichtet. Die Matrix verknüpft je Feld die \emph{geforderten Ansätze} mit dem
beobachtbaren Abdeckungsgrad sowie geeigneten Evidenzf




% -------------------------
% CHAPTER 6: Coverage-Analyse entlang ISO/PAS 8800
% -------------------------
\chapter{Coverage-Analyse entlang ISO/PAS 8800}\label{chap:coverage}

\section{Vorgehen der Analyse}\label{sec:analyse-vorgehen}
Die Analyse fokussiert nicht die Wirksamkeit einzelner Techniken, sondern \emph{welche Ansätze zur Absicherung gefordert werden}
(u.\,a.\ aus ISO/PAS~8800 sowie konzeptionellen Leitlinien) und \emph{in welchem Umfang Abdeckung erkennbar ist}.
Als Anker dienen die in Kapitel~\ref{chap:anforderungen} definierten \emph{Anforderungsfelder}. 
Für jedes Feld wird (i) der Bezug zu Forderungsquellen notiert, (ii) typische Evidenz/Artefakte skizziert, 
(iii) ein Abdeckungsgrad bewertet und (iv) eine kurze Maßnahme abgeleitet. 

\section{Anforderungs- und Coverage-Matrix}\label{sec:coverage}

% ===== Bewertungsrubrik & Ampel (neu) =====
\definecolor{ampelgruen}{RGB}{0,153,0}
\definecolor{ampelgelb}{RGB}{232,168,0}
\definecolor{ampelrot}{RGB}{194,0,0}

\newcommand{\ampelDot}[1]{%
  \begingroup
  \def\c{#1}%
  \ifx\c\relax \fbox{?}\else
  \ifx\c\empty \fbox{?}\else
  \ifx\c G {\color{ampelgruen}\Large$\bullet$}\fi
  \ifx\c Y {\color{ampelgelb}\Large$\bullet$}\fi
  \ifx\c R {\color{ampelrot}\Large$\bullet$}\fi
  \fi\fi\endgroup
}

% 0–3 Skala mit Ampelpunkt (0=R, 1=Y, 2=Y, 3=G)
\newcommand{\score}[1]{%
  \ifnum#1=0 \ampelDot{R}\,#1\fi%
  \ifnum#1=1 \ampelDot{Y}\,#1\fi%
  \ifnum#1=2 \ampelDot{Y}\,#1\fi%
  \ifnum#1=3 \ampelDot{G}\,#1\fi%
}

\newcommand{\ampelLegende}{%
\textbf{Bewertungsrubrik:} \score{0} = nicht adressiert,\;
\score{1} = partiell,\;
\score{2} = weitgehend,\;
\score{3} = adressiert \& nachweisbar.%
}

% Marker für Fokus (Frontloading/Inline)
\newcommand{\F}{\textsuperscript{F}} % Frontloading-bezogen
\newcommand{\I}{\textsuperscript{I}} % Inline-bezogen

\begin{center} % optional fürs Zentrieren
\begin{longtable}{p{0.20\textwidth} p{0.21\textwidth} p{0.25\textwidth} p{0.07\textwidth} p{0.27\textwidth}}
\caption{Coverage-Analyse (ISO/PAS~8800 \textrightarrow\ Anforderungsfelder)}\label{tab:coverage}\\

\toprule
\textbf{Anforderungsfeld} & \textbf{ISO/PAS~8800 Bezug} & \textbf{Erwartete Artefakte / Nachweise} & \textbf{Abde- ckung} & \textbf{Kommentar / Lücke / Maßnahmen} \\
\midrule
\endfirsthead

\multicolumn{5}{l}{\small Fortsetzung von \tablename~\thetable: Coverage-Analyse (ISO/PAS~8800 \textrightarrow\ Anforderungsfelder)}\\
\toprule
\textbf{Anforderungsfeld} & \textbf{ISO/PAS~8800 Bezug} & \textbf{Erwartete Artefakte / Nachweise} & \textbf{Abde- ckung} & \textbf{Kommentar / Lücke / Maßnahmen} \\
\midrule
\endhead

\midrule
\multicolumn{5}{r}{\small Fortsetzung auf der nächsten Seite}\\
\endfoot

\bottomrule
\endlastfoot

% --------- TABELLENINHALT (unverändert) ---------
\textbf{ODD-/Szenario- abdeckung} & ODD, Annahmen, szenarienbasiertes Testen & ODD-Dokument, Szenariokatalog, segmentiertes Leistungs-/Unsicherheits\-reporting, Re-Assessment bei ODD-Änderung & \score{2} & ODD-Schnitte definiert. \textbf{Maßn.:} verbindliches ODD-Template; Pflichtmetriken (z.B. Pd/Pfa) je Segment; Unknown/Abstain-Kriterien; Regression bei ODD-Änderung. \\
\midrule
\textbf{Daten-/Label-Governance}\F & Daten-Management, Dokumentation, Freigabeprozesse & Datasheets/Model Cards, Label-Guidelines, versionierte Datenstände (Signaturen) & \score{2} & Prozesse vorhanden. \textbf{Maßn.:} Freigabe-Checkliste, Pipeline-Hashing/Signierung, Datenänderungslog. \\
\midrule
\textbf{Robustheit \& Verteilungssprünge} & Robustheitsprinzi- pien, Stresstests & Stress-/Sensitivitätstests (z.\,B.\ RFI/Clutter/Noise), Degradationskurven & \score{1} & Systematische Stresstests lückenhaft; RFI/Multipath selten. \textbf{Maßn.:} Testbatterie inkl.\ Worst-Case, Katalog pro Störart, Akzeptanzkriterien. \\
\midrule
\textbf{Konfidenzkali- brierung \& Unsicherheit} & Unsicherheitsbe- handlung, Freigabekriterien & Kalibrationsberichte je ODD-Segment, dokumentierte Thresholds/Abstain-Regeln & \score{2} & Konfidenzmetriken vorhanden, Kalibrierung nicht konsistent je Segment. \textbf{Maßn.:} Kalibrierung verpflichtend pro Segment; Bootstrap-Reporting. \\
\midrule
\textbf{Toolchain-Assurance \& Reproduzierbarkeit}\F & Tool-/Prozesskontrolle, ggf.\ Toolqualifizierung & deterministische Builds, Lockfiles, Artefakt-Signaturen, reproduzierbare Trainings-/Eval-Läufe, TQL-Bewertung & \score{2} & Repro besser, formale Toolqualifizierung teils offen. \textbf{Maßn.:} TQL-Entscheidungsmatrix, Build-Manifeste. \\
\midrule
\textbf{Monitoring/ Degradation (Betrieb)}\I & Betriebsüberwa- chung, sichere Degradation & Laufzeitmetriken (Leistung/Kalibrierung/OOD-Rate), Drift-Detektion, Alarm-/Fallback-Policy & \score{1} & Degradationspfade wenig erprobt. \textbf{Maßn.:} OOD-/Drift-Tracker, Degradation-Playbook mit Tests. \\
\midrule
\textbf{Traceability/ Versionierung}\F & Rückverfolgbarkeit, Konfigurationsmgmt. & End-to-end Links Anf.\,$\rightarrow$\,Daten\,$\rightarrow$\,Modell\,$\rightarrow$\, Tests\,$\rightarrow$\,Release; Änderungsprotokolle & \score{2} & Links vorhanden, Lücken bei Daten\,$\leftrightarrow$\,Modell. \textbf{Maßn.:} Traceability-Matrix, verbindliche Artefakt-IDs. \\
\midrule
\textbf{XAI \& Auditierbarkeit}\F & nachvollziehbare Evidenz & zweckgebundene Erklärungen, reproduzierbare Notebooks/Skripte, Review-Artefakte & \score{2} & Heterogene XAI-Methoden, Repro schwankt. \textbf{Maßn.:} XAI-Playbook, Seed/Env-Fixierung, Audit-Pipeline. \\
\midrule
\textbf{Änderungen/ Freigabe}\F & Änderungs- und Freigabekontrolle & Re-Training-Trigger, Regression-Testplan, Freigabeprotokoll, Safety-Case-Update & \score{2} & Trigger uneinheitlich. \textbf{Maßn.:} feste Trigger (Daten-/Shift/Performance), definierter Regression-Scope. \\
\midrule
\textbf{Schnittstellen/ HMI (Inline)}\I & Informationsüber- gabe, Interpretierbarkeit & Qualitätsflags, Konfidenz, Unknown/Abstain, HMI-Richtlinien (OOD-Hinweise) & \score{2} & Flags vorhanden, Darstellung uneinheitlich. \textbf{Maßn.:} HMI-Styleguide, Operator-Validierung. \\
\midrule
\textbf{Frontloading-spezifisch}\F & Evidenzgenerierung, Transfer in Safety Case & Audit-Trails (Notebooks/Datenstände), Evidenzpakete, Review mit Safety-Team & \score{2} & Evidenztransfer inkonsistent. \textbf{Maßn.:} Evidenz-Template, Freigabe-Checklist, gemeinsame Reviews. \\

\end{longtable}
\end{center}

\vspace{0.5em}
\ampelLegende\quad \F\,= Frontloading-bezogen,\; \I\,= Inline-bezogen.


\section{Diskussion der Abdeckung}\label{sec:abdeckung-diskussion}
Die Coverage-Analyse zeigt ein differenziertes Bild über alle Anforderungsfelder hinweg. Insgesamt ist eine \textbf{weitgehende Abdeckung}
in prozessualen und evidenzbezogenen Feldern erkennbar (z.\,B.\ Daten-/Label-Governance, Toolchain-Assurance), während
\textbf{leistungs- und betriebskritische Felder} (Robustheit \& Verteilungssprünge, Monitoring/Degradation) noch \textbf{Lücken}
aufweisen. Die wichtigsten Beobachtungen sind:

\paragraph{(1) ODD-/Szenarioabdeckung: weitgehend vorhanden, Unsicherheits-Reporting inkonsistent.}
Die Definition der ODD und die szenarienbasierte Leistungsbewertung sind konzeptionell verankert. \textit{Schwachpunkte} liegen in der
\textit{einheitlichen Berichterstattung von Unsicherheit} pro Segment (Kalibrierung, Konfidenzbänder) sowie in \textit{klaren Kriterien für Unknown/Abstain}.
Eine verbindliche ODD-Vorlage und Pflichtmetriken (z.\,B.\ Pd/Pfa je Segment) würden die Nachvollziehbarkeit erhöhen und
Re-Assessments bei ODD-Änderungen operationalisieren.

\paragraph{(2) Daten-/Label-Governance: organisatorisch gut adressierbar, in der Tiefe heterogen.}
Prozesse zu Herkunft, Qualität, Freigabe und Dokumentation sind \textbf{umsetzbar und auditierbar}. In der Praxis variieren jedoch
\textbf{Label-Qualität}, Versionierungstiefe und Freigabe-Checklisten. Nicht ausreichend \textbf{Standardisierte Datasheets/Model Cards} und \textbf{signierte Datenstände}
reduzieren Reproduzierbarkeits- und Integritätsrisiken insbesondere im Frontloading.

\paragraph{(3) Robustheit \& Shift: größter Handlungsbedarf.}
Systematische \textbf{Stress- und Shift-Tests} sind uneinheitlich etabliert; \textbf{RFI/Multipath/Clutter}-Spezifika werden häufig nur punktuell adressiert.
Es fehlt eine \textbf{Testbatterie} mit Worst-Case-Definitionen und akzeptierten \textbf{Degradationskurven}. Solange diese Evidenzen nicht belastbar vorliegen,
bleibt die Bewertung „sicher genug“ \textbf{konditional}.

\paragraph{(4) Konfidenzkalibrierung \& Unsicherheit: verfügbar, aber nicht durchgängig segmentiert.}
Konfidenzmetriken liegen vor; \textbf{Kalibrierung je ODD-Segment} und \textbf{Threshold-/Abstain-Regeln} sind jedoch nicht konsistent dokumentiert.
Eine verpflichtende \textbf{segmentierte Kalibrierung} stärkt die Freigabekriterien und erleichtert die Auditierung.

\paragraph{(5) Toolchain-Assurance \& Traceability: solide Basis, Formalisierung ausbaufähig.}
Reproduzierbare Builds, Lockfiles, Artefakt-Signaturen und End-to-End-Traceability sind vielfach vorhanden, jedoch \textbf{nicht lückenlos}.
\textbf{TQL-Entscheidungsmatrizen} (wo relevant), verbindliche \textbf{Artefakt-IDs} und \textbf{Pflicht-Change-Logs} schließen typische Nachweislücken.

\paragraph{(6) Monitoring/Degradation im Betrieb: reif für Leitplanken, noch nicht Standard.}
Laufzeitmetriken und Drift-/OOD-Detektoren werden \textbf{rudimentär} genutzt; \textbf{klare Degradationspfade} (Alarmierung, Fallback) fehlen oft.
Ein \textbf{Degradation-Playbook} mit Testfällen und Telemetrie-Set ist Voraussetzung für verlässliche Betriebsevidenz, insbesondere für Inline-Funktionen.

\medskip
\noindent\textbf{Frontloading vs.\ Inline.} Die Abdeckung ist im \textbf{Frontloading} (Daten-/Modellevidenz, Auditierbarkeit, Tool-/Prozesskontrolle)
\textbf{deutlich stärker} ausgeprägt als für \textbf{Inline-Funktionen}, die \textbf{Echtzeit-Monitoring und sichere Degradation} verlangen.
Die Matrix legt nahe, dass \textbf{Evidenzgenerierung und Nachweisführung} im Frontloading heute schon \textbf{plausibel strukturierbar} sind,
während betriebsnahe Nachweise (Robustheit unter Störung, Laufzeit-Drift, HMI-Integrationsqualität) den Hauptteil des verbleibenden Risikos tragen.

\medskip
\noindent\textbf{Implikation für die Leitfrage.} Aus Meta-Sicht entsteht kein binäres „Ja/Nein“, sondern eine \textbf{konditionale Einschätzung}:
Die \textbf{geforderten Ansätze} sind in mehreren Feldern \textbf{adressierbar} und \textbf{teilweise abgedeckt}; zugleich bestehen \textbf{wesentliche Evidenzlücken}
bei Robustheit/Shift und im Betrieb (Monitoring/Degradation). Vertrauen ist dort begründbar, \textbf{wo} ODD klar begrenzt, Evidenzen \textbf{segmentiert}
geführt und Degradationspfade definiert sind; darüber hinaus ist \textbf{gezielte Weiterentwicklung} erforderlich.

\medskip
\noindent\textbf{Überleitung.} Kapitel~\ref{chap:uebertragbarkeit} bewertet im Anschluss, inwieweit die geforderten Ansätze
\textbf{unter den Rahmenbedingungen von Radarsystemen} (Frontloading und Inline) \textbf{plausibel übertragbar} sind, und leitet daraus
\textbf{priorisierte Empfehlungen} ab.






% -------------------------
% CHAPTER 7: Bewertung der Übertragbarkeit
% -------------------------
\chapter{Bewertung der Übertragbarkeit}\label{chap:uebertragbarkeit}

Dieses Kapitel bewertet, inwieweit die in Kapitel~\ref{chap:anforderungen} definierten Anforderungsfelder und die in
Kapitel~\ref{chap:coverage} verdichtete Coverage-Analyse \emph{plausibel} unter den Rahmenbedingungen von
Radarsystemen adressiert werden können. Die Einordnung bleibt produktneutral und folgt dem Meta-Ansatz:
Im Fokus steht nicht die Wirksamkeit einzelner Techniken, sondern ob die \emph{geforderten Ansätze} mit
heutigen Praktiken belastbar abgedeckt werden.

\section{Abdeckung der Anforderungen des Beispiel Modells}

\section{Frontloading-Use-Cases (exemplarisch)}\label{sec:uebertragbarkeit-frontloading}
Frontloading-Anwendungen (außerhalb der Inline-Funktion; z.\,B.\ Datenaufbereitung, Musteranalyse, Validierung
und Evidenzerstellung) zeigen eine hohe Anschlussfähigkeit an die geforderten Ansätze in den Feldern
\emph{Daten-/Label-Governance}, \emph{Toolchain-Assurance \& Reproduzierbarkeit}, \emph{Traceability/Versionierung}
und \emph{Konfidenzkalibrierung \& Unsicherheit}.
\begin{itemize}
  \item \textbf{Stärken:} Reproduzierbare Pipelines (Seeds/Umgebungen), versionierte und signierte Artefakte,
  auditierbare Notebooks/Skripte und definierte Review-Prozesse ermöglichen belastbare Evidenzpakete
  für den Safety Case.
  \item \textbf{Schwachpunkte:} Die \emph{Tiefe der Robustheitsevidenz} gegenüber Verteilungssprüngen ist oft uneinheitlich;
  Stress-/Shift-Testkataloge und Degradationskurven sind nicht immer standardisiert.
  \item \textbf{Implikation:} Frontloading ist ein tragfähiges Vehikel zur \emph{Evidenzgenerierung}, sofern ODD-Schnitte,
  segmentierte Kalibrierung und Freigabekriterien verbindlich festgelegt und konsequent in Evidenzpakete überführt werden.
\end{itemize}

\paragraph{Warum der Schwerpunkt auf Frontloading.}
In klassischen Radarsystemen werden Mustererkennung und Schwellwerte oft fest verdrahtet: Schwellen entstehen aus der
Auswertung historischer Daten und werden anschließend deterministisch implementiert. Diese Musterfindung ist heute
Gegenstand der Datenwissenschaft, deren leistungsfähigste Werkzeuge (u.\,a.) KI-Methoden sind. Wird KI \emph{vorlaufend}
(offboard) zur Datenanalyse, Musterentdeckung und Hypothesenbildung eingesetzt – und das operative Verhalten anschließend
\emph{klassisch} (deterministisch) umgesetzt und nach DO-178C-Systemlogik verifiziert –, bleibt die KI vom Flugzeugbetrieb
entkoppelt und entspricht der EASA-Logik eines Level-1-Einsatzes (unterstützend, ohne eigenständige Autorität).
Der assurance­relevante Vorteil: Die KI dient als Erkenntnisinstrument, während die produktive Funktion durch konventionelle,
prüfbare Artefakte (z.\,B.\ feste Regeln, nachvollziehbare Schwellen) realisiert wird. Dadurch verlagern sich KI-spezifische
Nachweise primär auf Reproduzierbarkeit der Analysen, Daten-/Label-Governance und die Traceability der Ableitung
(Daten\,$\rightarrow$\,Muster\,$\rightarrow$\,Spezifikation), statt auf laufzeitkritische Nachweise im Bordbetrieb.

\paragraph{Hardware-Perspektive: Frontloading vs.\ Inline.}
Im \emph{Frontloading} liegen die Rechenlast und Speicherbedarfe offboard (Skalierung via CPU/GPU/Acceleratoren, kein Echtzeitzwang). 
Das reduziert hardwarebedingte Zulassungsrisiken: deterministische Umgebungen (Container/VM, Seeds, Lockfiles) genügen für 
Reproduzierbarkeit; Bordhardware bleibt unverändert, da die operative Logik \emph{deterministisch} implementiert wird. 
Im \emph{Inline}-Betrieb treten zusätzliche Anforderungen auf: harte Latenz- und Verfügbarkeitsbudgets, deterministisches 
Laufzeitverhalten unter Worst-Case-Last, begrenzte Ressourcen (Thermik, Energie, Speicher), potenzieller Bedarf an 
beschleunigender Hardware (z.\,B.\ GPU/NPU) mit Integrations- und Lebenszyklusrisiken (Obsoleszenz), sowie Implikationen für 
Nachweis und Qualifizierung der Hardware/Toolkette. Diese Zwänge erhöhen die Komplexität von Integration, Test und Zulassung 
gegenüber einer reinen Frontloading-Nutzung.



\section{Inline-Funktionen (Überblick)}\label{sec:uebertragbarkeit-inline}
Für Inline-Funktionen (laufzeitwirksame Nutzung im Betrieb) treten betriebsnahe Forderungen in den Vordergrund:
\emph{Monitoring/Degradation}, robuste Performance unter Störung sowie \emph{schnittstellen- und HMI-seitige}
Informationsqualität (Qualitätsflags, Konfidenz, Unknown/Abstain).
\begin{itemize}
  \item \textbf{Stärken:} Segmentierte Metriken und ODD-Begrenzung sind konzeptionell vorhanden; erste OOD-Indikatoren und
  Konfidenzberichte werden genutzt.
  \item \textbf{Schwachpunkte:} Es fehlen häufig integrierte \emph{Betriebsstrategien} mit eindeutigem Telemetrie-Set,
  Alarmierungslogik und getesteten Degradationspfaden; systematische RFI/Clutter-Stresstests sind lückenhaft.
  \item \textbf{Implikation:} Die Übertragbarkeit ist \emph{konditional}: klare ODD-Grenzen, definierte Abbruch-/Fallback-Regeln
  und evidenzgestützte Kalibrierung im Zielbetrieb sind Voraussetzung.
\end{itemize}

\section{Zusammenführende Bewertung}\label{sec:uebertragbarkeit-zusammen}
In Summe ist die Übertragbarkeit der \emph{geforderten Ansätze} für Frontloading-Szenarien heute \emph{plausibel} und mit
überschaubaren Ergänzungen (z.\,B.\ Standardisierung von Evidenzformaten, konsistente Kalibrierung je ODD-Segment)
erreichbar. Für Inline-Funktionen bestehen weiterhin \emph{kritische Evidenzbedarfe} bei \emph{Robustheit \& Shift}
sowie \emph{Monitoring/Degradation}. Die leitende Frage „\emph{ist KI sicher genug?}“ lässt sich daher nicht binär,
sondern \emph{konditional} beantworten: Dort, wo ODD eindeutig begrenzt, segmentierte Evidenz vorliegt und
Degradationspfade geprüft sind, ist Vertrauen begründbar; darüber hinaus bleibt gezielte Weiterentwicklung erforderlich.

\section{Brücke zu Empfehlungen}\label{sec:uebertragbarkeit-bruecke}
Die nachfolgenden Empfehlungen (Kapitel~\ref{chap:empfehlungen}) priorisieren Maßnahmen, die die in Kapitel~\ref{chap:coverage}
sichtbar gewordenen Lücken (insbesondere Robustheit/Shift und Betriebsevidenz) adressieren und die Operationalisierung
der geforderten Ansätze in Richtung auditierbarer Evidenzketten beschleunigen.

